% !TEX TS-program = lualatex
Диаграмма \refimage{game-services-interaction} описывает жизненный цикл взаимодействия пользователя с веб-интерфейсом системы игры в нарды. В архитектуре участвуют следующие компоненты: пользователь, браузер (клиентская часть), сервер игры (Game Server), сервер моделей (Model Server) и хранилище состояния (State Repository). Взаимодействие организовано по HTTP-протоколу с использованием REST API.

\image[width=0.9\textwidth]{Диаграмма взаимодействия пользователя с сервисами игры}{game-services-interaction}{game-services-interaction}

Процесс начинается с инициализации сессии. Пользователь открывает сайт, в результате чего браузер отправляет HTTP-запрос GET / на сервер игры. Сервер проверяет наличие открытой сессии. Если сессия не обнаружена, генерируется уникальный идентификатор пользователя (UUID), создаётся новая сессия, и клиент получает HTML-страницу с установленной cookie session[uid]. Браузер загружает необходимые статические ресурсы (JavaScript, CSS) и отображает пользователю игровое поле.

Первым действием пользователя в игровой сессии является нажатие кнопки <<Бросить кости>>. Браузер отправляет запрос POST /roll\_dice на сервер игры. Сервер извлекает UID из сессии, самостоятельно генерирует значения костей и возвращает их обратно клиенту. Браузер обновляет интерфейс и показывает пользователю результат броска.

После этого пользователь выбирает один из допустимых ходов и инициирует следующий запрос: POST / с телом, содержащим текущие значения костей, состояние доски и выбранное действие. Сервер снова извлекает UID из сессионной cookie и пересылает запрос на Model Server по маршруту /make\_turn, включая UID в заголовках запроса.

Model Server обращается к State Repository для получения текущего состояния игры, связанного с данным пользователем. В случае, если это первый ход, создаётся новая запись состояния. Затем сервер обрабатывает действие пользователя, применяет модель TD-Gammon для вычисления хода искусственного интеллекта, обновляет игровое состояние и сохраняет его обратно в репозиторий. После завершения вычислений Model Server возвращает серверу игры данные: ответный ход AI, новое состояние доски и статус игры (продолжается или завершена).

Сервер игры передаёт полученные данные браузеру. Интерфейс обновляет игровую доску и показывает пользователю сделанный AI-ход. В зависимости от статуса игры, либо отображается итог (если игра завершена), либо интерфейс остаётся в ожидании следующего хода пользователя.

В системе предусмотрены механизмы обработки исключительных ситуаций. Если сессия пользователя отсутствует или истекла, запрос к серверу будет отклонён с кодом 401 Unauthorized, а клиент предложит пользователю перезагрузить страницу. В случае, если Model Server недоступен (например, из-за таймаута или внутренней ошибки), сервер игры вернёт ошибку 503 Service Unavailable, и браузер отобразит соответствующее сообщение.
